\chapter{Testing and Verification}
\label{ch:methodology}

\section{Verification Approach}

Verification was carried out at two complementary levels: RTL-level
simulation before synthesis, and hardware-level functional verification
on the physical FPGA after programming. This two-level approach is
standard in digital design practice~\citep{sze2017}. RTL simulation is
fast and exposes internal signals for inspection, but it cannot on its
own guarantee that timing closure and synthesis artefacts have not
introduced discrepancies between the simulated and deployed design;
hardware verification on the device closes that gap by exercising the
complete SoC end-to-end. A discrepancy at either level is sufficient to
fail a case, and every recorded cycle count in Chapter~\ref{ch:results}
was taken from a run that passed both.

Within the simulation level the testbench suite was deliberately
structured to combine the three verification styles used in
industry-grade digital flows: \emph{directed tests} for targeted
functional cases with hand-computed golden outputs, \emph{concurrent
SystemVerilog Assertions (SVA)} for always-on protocol and liveness
monitoring, and \emph{constrained-random stimulus with functional
coverage and a scoreboard} for exploring the input space beyond what
directed tests reach. Four SystemVerilog testbenches cover the design:
three unit testbenches under \texttt{tb/unit/} that target the
accelerator modules in isolation, and one integration testbench under
\texttt{tb/integ/} that exercises the full synthesisable SoC with
firmware loaded into a behavioural RAM model.

\subsection*{Unit Testbench: Datapath}

\texttt{accel\_datapath\_tb.sv} exercises the arithmetic datapath in
isolation. It contains six directed tests covering $B_\mathrm{Seg}$
load, MAC accumulation, C-tile clear, ReLU activation (positive, zero,
and negative inputs), signed MAC with negative operands, and a small
hand-computed $2 \times 2$ SpMM. On top of the directed tests the
testbench runs a constrained-random MAC-scoreboard test: a transaction
class with \texttt{rand} row and operand fields is randomised for each
of twenty MAC operations per iteration, and a parallel software model
of the C-tile is updated in lock-step with the DUT. At the end of each
iteration every accessible C-tile location is read back and compared
against the software reference; any mismatch is reported with the
exact $(\mathrm{row}, \mathrm{col})$ index and stops the test. Five
concurrent SVA properties run for the entire simulation: ReLU output
non-negativity when enabled, writeback pass-through when disabled,
one-hot datapath command exclusivity (clear, $B_\mathrm{Seg}$ write,
MAC, and writeback may not fire simultaneously), MAC row index in
range, and $B_\mathrm{Seg}$ index in range. A functional covergroup
records MAC row band (low / mid / high) crossed with operand sign
(positive / zero / negative) so that the random test has a visible
coverage metric for how much of the tile and operand space has been
exercised.

\subsection*{Unit Testbench: Control FSM}

\texttt{accel\_ctrl\_tb.sv} verifies control-path behaviour by
instantiating the control module together with a synchronous RAM model
with the same one-cycle read latency as the deployed BRAM, plus a
datapath stub. Six concurrent SVA properties monitor the FSM's external
contract for the whole simulation: single-port RAM exclusivity (read
and write strobes may not be simultaneously high), 32-bit address
alignment on every RAM access, quiescence of the RAM interface during
reset, mutual exclusion of the \texttt{BUSY} and \texttt{DONE} status
bits, FSM liveness (every \texttt{BUSY} rise must be followed by a
\texttt{BUSY} fall within one million cycles), and one-hot exclusivity
of the four datapath command outputs (\texttt{dp\_clear\_en},
\texttt{dp\_bseg\_we}, \texttt{dp\_mac\_en}, \texttt{dp\_wb\_en}). A
directed smoke test drives the FSM through a full start--busy--done
cycle and checks that the IRQ line asserts correctly when
\texttt{irq\_en} is set.

\subsection*{Unit Testbench: Accelerator Top}

\texttt{accel\_top\_tb.sv} is the most comprehensive unit testbench and
closes the loop across the MMIO front-end, the FSM, and the datapath.
Stimulus is structured into two classes following the
\emph{transaction / driver} pattern common in UVM-style flows but
implemented in plain SystemVerilog: \texttt{SpmmConfig} is a
transaction object with \texttt{rand} fields for $M$, $N$, $K$, and
\texttt{relu\_en} and an embedded covergroup over the relevant
combinations; \texttt{accel\_drivers} holds a virtual handle to the
DUT's MMIO interface and implements the \texttt{mmio\_read},
\texttt{mmio\_write}, and \texttt{wait\_done} protocol tasks so that
every test uses the same low-level interface sequencing. Five directed
tests verify MMIO register read/write coverage, the empty-matrix
($\mathrm{NNZ} = 0$) case, a hand-computed $2 \times 2 \times 8$ SpMM,
ReLU clamping of a negative partial sum, and the BUSY-mirror LED. A
sixth constrained-random test randomises twenty \texttt{SpmmConfig}
objects, runs each one through the full accelerator end-to-end, reads
back the output C matrix, and verifies element-by-element that it
matches the trivially computable all-zero golden output; after each
iteration the covergroup sampling is reported so that coverage
convergence is visible during simulation. Five concurrent SVA
properties guard invariants that must hold across every one of these
tests: single-port RAM exclusivity, RAM address alignment, reset
quiescence, \texttt{BUSY}/\texttt{DONE} status mutex, and top-level FSM
liveness. Different simulation seeds (\texttt{+SEED=N}) exercise
different configuration combinations; coverage-driven regressions
simply increase the seed count until every bin of the covergroup is
hit.

\subsection*{Integration Testbench: Full SoC}

\texttt{soc\_tb.sv} instantiates PicoRV32 together with the accelerator
and a 64\,KB behavioural RAM, loads the compiled firmware image into
the RAM via \texttt{\$readmemh}, releases reset, and lets the firmware
execute its full benchmark sequence until the CPU writes a sentinel
value to a pre-agreed result address. This closes the loop between the
software control flow described in Chapter~\ref{ch:firmware} and the
hardware being verified: any firmware change that breaks the MMIO
contract, any RTL change that breaks the firmware's assumed timing, or
any memory-map misconfiguration fails immediately in simulation before
being committed to a synthesis run. Five concurrent SVA properties
monitor the shared memory fabric: address-decode exclusivity between
the RAM region and the accelerator MMIO region, bounded CPU handshake
progress (\texttt{mem\_ready} must rise within 1024 cycles of
\texttt{mem\_valid}), accelerator RAM 32-bit address alignment,
accelerator RAM single-port exclusivity, and an IRQ-epoch guard that
ensures the accelerator's IRQ line can only rise after the accelerator
has actually been started.

\subsection*{Shared Infrastructure}

All four testbenches share the support files under \texttt{tb/common/}:
\texttt{clk\_reset.sv} provides the clock generator and reset task,
\texttt{tb\_pkg.sv} provides plusarg helpers
(\texttt{get\_plusarg\_int}, \texttt{has\_plusarg}) so every testbench
accepts \texttt{+SEED} and \texttt{+WAVES} uniformly,
\texttt{tb\_macros.svh} defines the line-number-tagged
\texttt{TB\_CHECK} and \texttt{TB\_FATAL} PASS/FAIL macros, and
\texttt{tb\_regmap.svh} centralises the MMIO register offsets and bit
positions so that they are declared once and \texttt{`include}d
everywhere, preventing silent drift between testbenches and RTL.

\subsection*{Hardware-Level Verification}

Hardware-level verification is embedded directly inside the firmware
benchmark harness described in Section~\ref{sec:data-prep}. For each
of the 12 cases and at each of the three accelerator stages, firmware
generates deterministic inputs from fixed seeds, computes the
quantised CPU reference using the \texttt{spmm\_cpu\_q8()} routine,
runs the accelerator on the same inputs, and invokes the
\texttt{array\_compare()} scoreboard to verify every element of the
accelerator output matches the software reference exactly. This
routine functions as a hardware-level golden-reference scoreboard: it
covers all $M \times N$ output elements on every run, and a single
incorrect element causes the case to fail and no performance number
to be recorded. The combination of RTL-level SVA monitoring,
randomised-scoreboard unit testbenches, and firmware-level
element-wise comparison on real silicon provides three independent
layers of functional checking before any cycle count is reported.

\section{Benchmark Suite Design}
\label{sec:bench-design}

The benchmark suite was designed to cover multiple independent axes of
performance variation so that the reasons behind speedup trends can be
identified analytically rather than anecdotally. Four axes were chosen:
matrix size ($M$, $K$), which governs the amortisation of fixed setup
overhead; dense output width $N$, which determines the number of tile
columns and the ratio of tile-management cycles to useful computation;
sparsity level, which sets NNZ and therefore the volume of useful
arithmetic per row traversal; and nonzero distribution pattern, which
exposes any sensitivity of the accelerator or CPU to the spatial layout
of the nonzeros. Table~\ref{tab:benchmark-inputs} lists all 12 cases
together with the axis of variation each one primarily addresses.

\begin{table}[htbp]
  \centering
  \scriptsize
  \setlength{\tabcolsep}{4pt}
  \begin{tabular}{llrp{2.6cm}p{5.4cm}}
  \toprule
  Test ID & $M \times K \times N$ & Sparsity & Pattern & Purpose \\
  \midrule
  T1  & $8 \times 8 \times 8$    & 75\% & Uniform &
    Small debug/sanity case; establishes minimum speedup in the suite. \\
  T2  & $16 \times 16 \times 8$  & 75\% & Uniform &
    Small size-scaling point; 4$\times$ more arithmetic than T1. \\
  T3  & $32 \times 32 \times 8$  & 75\% & Uniform &
    Medium size-scaling point. \\
  T4  & $64 \times 64 \times 8$  & 75\% & Uniform &
    Largest square case; maximum $M$ for the current $M_{\max} = 64$. \\
  T5  & $64 \times 64 \times 16$ & 75\% & Uniform &
    Tests effect of doubling dense width relative to T4. \\
  T6  & $64 \times 64 \times 32$ & 75\% & Uniform &
    Tests effect of quadrupling dense width; 4 tile columns. \\
  T7  & $32 \times 32 \times 32$ & 50\% & Uniform &
    Low sparsity (dense); more nonzeros per row than T8/T9. \\
  T8  & $32 \times 32 \times 32$ & 75\% & Uniform &
    Mid-sparsity reference; matches T10/T11 except pattern. \\
  T9  & $32 \times 32 \times 32$ & 90\% & Uniform &
    Very sparse; exposes overhead-to-arithmetic ratio. \\
  T10 & $32 \times 32 \times 32$ & 75\% & Row-skewed NNZ &
    Same NNZ as T8; row imbalance sensitivity. \\
  T11 & $32 \times 32 \times 32$ & 75\% & Clustered NNZ &
    Same NNZ as T8; block locality sensitivity. \\
  T12 & $64 \times 64 \times 32$ & 90\% & Uniform &
    Large sparse stress case; wide output and few nonzeros. \\
  \bottomrule
  \end{tabular}
  \caption{Complete benchmark input set used across all recorded
           evaluation stages.}
  \label{tab:benchmark-inputs}
\end{table}

\section{Cycle-Count Definitions and Speedup}

CPU cycle counts are measured around the timed software SpMM function
only. The measurement brackets the outer row loop, all CSR index reads
(\texttt{rowptr}, \texttt{colidx}), the inner dense-column MAC loop, and
all BRAM reads and writes to $B$ and $C$; matrix generation,
quantisation, and output-buffer clearing happen outside the bracket and
are deliberately excluded. Accelerator cycle counts are measured around
the complete \texttt{accel\_run\_spmm\_int8()} call and therefore include
the \texttt{clear\_done} write, the ten register writes of
\texttt{accel\_configure()} (roughly 20 cycles of overhead), the
\texttt{start} pulse, the complete hardware execution time, and the
polling loop that follows (about five cycles per iteration, typically no
more than ten iterations by the time \texttt{done} is observed). This
is an end-to-end offload latency from the CPU's perspective, which is
the correct metric for characterising user-visible benefit, but it is
explicitly not a measurement of pure datapath throughput.

Speedup is defined as
\[
  \text{Speedup} = \frac{T_{\mathrm{CPU}}}{T_{\mathrm{ACCEL}}},
\]
with both times measured in CPU clock cycles on the same 50\,MHz clock.
For the RISC-V comparisons $T_{\mathrm{CPU}}$ is the timed PicoRV32
software result (full INT32 for the baseline and DSP stages, INT8 for the
INT8 stage); for the Cortex-M0 comparison $T_{\mathrm{CPU}}$ is the
number of Cortex-M0 cycles computed from an equivalent instruction-count
model for the same benchmark case.

\section{Fairness, Limitations, and Threats to Validity}

The comparison is fair in several important respects. The CPU and
accelerator operate on the same benchmark inputs generated from the same
fixed seeds, they use the same shared on-chip RAM system at the same
50\,MHz clock, the software reference uses the M-extension \texttt{MUL}
instruction for its inner multiplies so that it is not artificially
slowed by software multiplication emulation, and the accelerator is
always validated against the software reference before its timing is
recorded. The Cortex-M0 comparison uses the same benchmark definitions
and the same accelerator cycle counts, with the CPU-side figures derived
from known Cortex-M0 instruction counts and the processor's published CPI
model (predominantly single-cycle instructions with 1--3 cycle load
latency). At the same time, the comparison differs in one important way:
the PicoRV32 software baseline runs on the very same in-order core that
controls the accelerator, with no cache, no speculation, and no SIMD. A
more aggressively optimised CPU implementation (one with
software-pipelined inner loops, block-reuse of $B$-rows, or explicit
prefetching) would narrow the gap. The goal of the baseline is
credibility and simplicity, not maximum CPU throughput.

Several limitations apply to all results presented in
Chapter~\ref{ch:results}. Each accelerator cycle count is a single
recorded value rather than a statistical distribution over multiple runs;
cycle counts for deterministic hardware are repeatable, but a
min/median/max summary would add an extra layer of confidence. The test
matrices are generated by a PRNG from fixed seeds rather than drawn from
real GNN or recommender-system datasets, so while they cover realistic
sparsity levels and pattern structures, performance on real workloads
may differ. All data lives in the 64\,KB on-chip BRAM; real SpMM
workloads typically span gigabytes and therefore incur DRAM latency,
and the results here reflect the best-case scenario of tightly coupled
on-chip memory. The accelerator cycle counts, as noted above, include
MMIO configuration and polling overhead, and for the smallest case (T1)
this overhead is a non-trivial fraction of the total. Finally, the
Quartus post-fit resource and timing reports were still being finalised
at the time of writing; the design has been verified to close timing at
50\,MHz, but the exact ALM, BRAM, and DSP counts across all three build
variants are pending.

These limitations define the scope of the claims that can be
responsibly made. The performance story (that INT8 packing delivers a
$2.18\times$ improvement over DSP MAC, which delivers a $1.32\times$
improvement over the baseline, and that memory bandwidth dominates the
cycle budget) is not undermined by any of them. But claims about
absolute performance at scale, DRAM bandwidth, or energy efficiency
would require a different experimental infrastructure than the one
developed here.
